\] Let \(\mathcal P_R\) be the real mean-zero trigonometric polynomials supported in \(\|v\|_\infty\leq R\). Its real dimension is \(4R(R+1)\).

\subsection{An exact finite theorem for arbitrary Fourier phases}\label{nondegeneracy:an-exact-finite-theorem-for-arbitrary-fourier-phases}

\textbf{Theorem 1 (the full radius-eight box).} For every nonzero \(g\in\mathcal P_8\), \[
                         \boxed{\Gamma_T(g)>0.}             \tag{ND.3}
\] In particular (ND.3) holds at every smaller integer radius. There is no restriction on the signs or phases of the Fourier coefficients. This proves that the second-order bracket is not a finite Fourier coboundary throughout this 288-dimensional real space. The assertion is not extrapolated to all radii.

\subsubsection{2.1 Exact polynomial equations}\label{nondegeneracy:exact-polynomial-equations}

Choose \[
 V_R=\{(a,b):0\leq a\leq R,\ -R\leq b\leq R,
                         \ a>0\text{ or }b>0\}.
\] Write, uniquely, \[
 g=\sum_{v\in V_R}(z_v e_v+\overline{z_v}e_{-v}),\qquad
 z_v\in\mathbb C,\qquad \|g\|_2^2=2\sum_v|z_v|^2.          \tag{ND.4}
\] There are \(n_R=2R(R+1)\) complex variables. The exact Fourier bracket is \[
 \widehat{\{g,\mathcal C_Tg\}/(4\pi)}(k)
       =-\pi\sum_{u+Tv=k}\det(u,Tv)\widehat g(u)\widehat g(v).
 \tag{ND.5}
\] Combining the coefficients along one of the two orientations of each signed orbit gives integer-coefficient quadratic polynomials \(Q_j(z,\bar z)\) with \[
                   \Gamma_T(g)=2\pi^2\sum_j|Q_j(z,\bar z)|^2. \tag{ND.6}
\] Thus (ND.3) is the assertion that all the equations \(Q_j=0\) have only the zero solution subject to the genuine complex-conjugation relation. The latter is not replaced by treating \(z\) and \(\bar z\) as independent complex variables.

\subsubsection{2.2 Why the orbit classification terminates and preserves phases}\label{nondegeneracy:why-the-orbit-classification-terminates-and-preserves-phases}

For \(v=(a,b)\) put \(\beta=b+\sqrt2 a\), \(\beta'=b-\sqrt2 a\). Then \[
 E(T^jv):=2(T^jv)_1^2+(T^jv)_2^2
       =\tfrac12(\beta^2\lambda^{2j}+\beta'^2\lambda^{-2j}). \tag{ND.7}
\] Both coefficients are positive for \(v\ne0\). This is strictly convex as a function of real \(j\), tends to infinity at both ends, and has at most two adjacent integer minimizers. Integer neighbour descent reaches the global minimum. At equal energies we use a fixed lexicographic rule, including the choice of a sign for the vector. Energy decreases in the non-tie steps, so the algorithm cannot loop indefinitely.

Crucially, the program records the orientation bit \(\epsilon\) in \(w=\epsilon T^jv\). A negative orientation contributes to the conjugate orbit, not a second independent copy of the same complex coefficient. Adjacent minimum ties, for example the orbit containing \((1,-2)\), are retained as explicit normalization tests. Omitting this bit can produce an incorrect arbitrary-phase variance even when all even-symbol tests pass.

The independent C++ constructor works in the quadratic ring with vector \(b+a\sqrt2\) and enumerates an orbit window \(-16\leq j\leq16\). It fails if a minimizer is at the boundary. Interior minimizers have the global meaning just established by (ND.7); the finite window is not an unproved coverage assumption. All 1,475 nonzero radius-eight equations agree exactly, term by term, with the independently constructed Python equations.

\subsubsection{2.3 A small proof language and its soundness}\label{nondegeneracy:a-small-proof-language-and-its-soundness}

